\documentclass{beamer}
\usepackage{graphicx}
\usepackage{xcolor}
\usepackage{tikz}

% 自定义颜色
\definecolor{purpleblue}{RGB}{80, 80, 180}
\definecolor{lightgray}{RGB}{230, 230, 230}
\definecolor{novelPurple}{RGB}{98, 54, 255} 
\definecolor{darkgray}{RGB}{50, 50, 50}

% 主题设置
\usetheme{Madrid}  
\usecolortheme{dolphin} 

% 设置字体（移除 fontspec，确保 pdfLaTeX 兼容）
\setbeamerfont{title}{size=\Large,series=\bfseries}
\setbeamerfont{frametitle}{size=\LARGE,series=\bfseries}
\setbeamercolor{frametitle}{fg=white, bg=novelPurple}
\setbeamercolor{title}{fg=white, bg=purpleblue}

% 封面信息
\title[SAT Hard Question]{\textbf{SAT Hard Question 1}}
\author{\textbf{Jack Zeng}}
\institute{\textcolor{white}{\textbf{Novel Prep}}}
\date{\today} 

\begin{document}

% --- Title Slide ---
\begin{frame}
    \titlepage
    \vfill
    \begin{flushright}
        \textcolor{purpleblue}{\Huge \textbf{Novel Prep}}
    \end{flushright}
\end{frame}


\begin{frame}
    \begin{tikzpicture}[remember picture, overlay]
        % 设置浅色渐变背景
        \shade[top color=softPurple, bottom color=softGray] 
            (current page.north west) rectangle (current page.south east);
        
        % 添加高斯模糊光晕
        \fill[white, opacity=0.3] (current page.center) circle (4cm);
        
        % 主标题
        \node[align=center] at (current page.center) {
            {\Huge\bfseries\textcolor{purpleblue}{Novel Prep}} \\[0.5cm]
            {\Large\itshape\textcolor{lightText}{Excellence in SAT Prep}}
        };

        % 添加柔和光点（点缀）
        \foreach \x in {1,2,3,4,5} {
            \fill[purpleblue, opacity=0.2] (rand*10-5, rand*7-3.5) circle (0.1+\x*0.05);
        }

        ;
    \end{tikzpicture}
\end{frame}


% Question 1
\begin{frame}{\textbf{Question: Greatest Possible Value }}
    \textbf{In the given expression,} \( a \) and \( c \) are positive constants.  
    If \( px + q \) is a factor of the expression:
    \[
    ax^2 + 176x + c = (px + q)(rx + s),
    \]
    where \( p \) and \( q \) are positive constants, what is the greatest possible value of \( ac \)?

    \vspace{0.5cm}
\end{frame}

% --- Answer Explanation 1 ---
\begin{frame}{\textbf{Answer Explanation: Maximum Value of \( ac \)}}
    \textbf{Correct Answer: \textcolor{novelPurple}{\( \mathbf{7744} \)}}  

    \vspace{0.3cm}
    \textbf{Solution:}
    \begin{itemize}
        \item We use the discriminant condition:
        \[
        \Delta = b^2 - 4ac > 0.
        \]
        \item Substituting \( b = 176 \):
        \[
        176^2 - 4ac > 0.
        \]
        \item Simplifying:
        \[
        4ac \leq 176^2.
        \]
        \item Solving for \( ac \):
        \[
        ac \leq 7744.
        \]
    \end{itemize}
    \textbf{Final Answer:} \( \boxed{7744} \).
\end{frame}



\begin{frame}{\textbf{Question: Finding the Greatest Possible Value of \( b \)}}
    In the given expression:
    \[
    34z^{18} + bz^9 + 30,
    \]
    \( b \) is a positive integer.  

    \vspace{0.3cm}
    If \( qz^9 + r \) is a factor of the expression, where \( q \) and \( r \) are positive integers, what is the greatest possible value of \( b \)?
\end{frame}

% --- Answer Explanation 2 ---
\begin{frame}{\textbf{Answer Explanation: Maximum Value of \( b \)}}
    \textbf{Correct Answer: \textcolor{novelPurple}{\( \mathbf{1021} \)}}  

    \vspace{0.3cm}
    \textbf{Solution:}
    \begin{itemize}
        \item Compute \( b \) using factor constraints:
        \[
        b = 34 \times 30 + 1.
        \]
        \item Performing the calculation:
        \[
        b = 1021.
        \]
    \end{itemize}
    \textbf{Final Answer:} \( \boxed{1021} \).
\end{frame}

% --- Question 3 ---
\begin{frame}{\textbf{Question: Finding the Value of \( a + c \)}}
    The expression:
    \[
    6x^4 + 17x^2 + 7
    \]
    can be rewritten as:
    \[
    (3x^2 + a)(2x^2 + b),
    \]
    where \( a \) and \( b \) are positive integers, or as:
    \[
    (3x^2 + c)(2x^2 + d),
    \]
    where \( c \) and \( d \) are positive non-integers.  

    \vspace{0.3cm}
    What is the value of \( a + c \)?
\end{frame}

\begin{frame}{Answer}
We are given:
\[
6x^4 + 17x^2 + 7
\]

We want to factor this as:
\[
(3x^2 + a)(2x^2 + b)
\]

Multiply:
\[
(3x^2 + a)(2x^2 + b) = 6x^4 + (2a + 3b)x^2 + ab
\]

Match coefficients:
\[
6x^4 + (2a + 3b)x^2 + ab = 6x^4 + 17x^2 + 7
\Rightarrow
\begin{cases}
ab = 7 \\
2a + 3b = 17
\end{cases}
\]




\textbf{Final Answer:} $\boxed{\dfrac{17}{2}}$
\end{frame}


% --- Question 4 ---
\begin{frame}{\textbf{Question: Finding the Smallest Possible Value of \( ab \)}}
    The given quadratic function:
    \[
    54x^4 + 219x^2 + 105
    \]
    has factors in the form:
    \[
    (k)(ax^2 + b)(cx^2 + d).
    \]
    If \( a, b, c, d, \) and \( k \) are integers, what is the smallest possible value of \( ab \)?
\end{frame}

\begin{frame}{Answer}
    \textbf{Final Answer:} $\boxed{14}$
\end{frame}


% --- Question 5 ---
\begin{frame}{\textbf{Question: Finding \( a - b - c \)}}
    The equation of a parabola is written in the form:
    \[
    y = ax^2 + bx + c,
    \]
    where \( a, b, \) and \( c \) are constants.  

    \vspace{0.3cm}
    When graphed in the \( xy \)-plane, the parabola has vertex \( (-1, 4) \) and does not intersect the \( x \)-axis.  

    \vspace{0.3cm}
    Which of the following could be the value of \( a - b - c \)?
    \textbf{Options:}
    \begin{enumerate}
        \item[A] 2
        \item[B] 4
        \item[C] -1
        \item[D] -6
    \end{enumerate}


    
\end{frame}




% --- Answer Explanation 5 ---
\begin{frame}{\textbf{Answer Explanation: Finding \( a - b - c \)}}
    \textbf{Correct Answer: \textcolor{novelPurple}{\( \mathbf{D} \)}}  

    \vspace{0.3cm}
    \textbf{Solution:}
    \begin{itemize}
        \item Since the parabola does not intersect the \( x \)-axis, the discriminant condition is:
        \[
        \Delta = b^2 - 4ac < 0.
        \]
        \item Solving the given system:
        \[
        \boxed{D}.
        \]
    \end{itemize}
\end{frame}

\begin{frame}{\textbf{Question: Maximum Value of \( b \)}}
    A quadratic function with vertex \( (1, 32) \) can be written as:
    \[
    -ax^2 + bx + c \quad (a, b, c \text{ are all positive integers}).
    \]
    What is the maximum value of \( b \)?
\end{frame}

% --- Answer Explanation 6 ---
\begin{frame}{\textbf{Answer Explanation: Maximum Value of \( b \)}}
    \textbf{Correct Answer: \textcolor{novelPurple}{\( \mathbf{62} \)}}  

    \vspace{0.3cm}
    \textbf{Solution:}
    \begin{itemize}
        \item Given:
        \[
        b = 2a \quad \text{and maximum } a + c = 32.
        \]
        \item Solving for \( b_{\text{max}} \):
        \[
        b_{\text{max}} = \boxed{62}.
        \]
    \end{itemize}
\end{frame}

% --- Question 7 ---
\begin{frame}{\textbf{Question: Converting Acceleration Units}}
    An object's speed is increasing at a rate of \( 7.7 \, \text{m/s}^2 \).  

    \vspace{0.3cm}
    What is this rate in miles per minute squared, rounded to the nearest tenth?  

    (Use \( 1 \, \text{mile} = 1,609 \, \text{meters} \).)

    \vspace{0.5cm}
    \textbf{Options:}
    \begin{enumerate}
        \item[A] 0.3
        \item[B] 17.2
        \item[C] 206.5
        \item[D] 209
    \end{enumerate}
\end{frame}

% --- Answer Explanation 7 ---
\begin{frame}{\textbf{Answer Explanation: Converting Acceleration Units}}
    \textbf{Correct Answer: \textcolor{novelPurple}{\( \mathbf{17.2} \)}}  

    \vspace{0.3cm}
    \textbf{Solution:}
    \begin{itemize}
        \item Convert units:
        \[
        7.7 \, \frac{\text{m}}{\text{s}^2} \times \left(\frac{60 \, \text{s}}{1 \, \text{min}}\right)^2 \times \frac{1 \, \text{mile}}{1609 \, \text{m}}.
        \]
        \item Simplify:
        \[
        7.7 \times \frac{3600}{1609}.
        \]
        \item Approximate:
        \[
        \approx 17.2 \, \frac{\text{miles}}{\text{min}^2}.
        \]
    \end{itemize}
    \textbf{Final Answer:} \( \boxed{17.2} \).
\end{frame}

% --- Question 8 ---
\begin{frame}{\textbf{Question: Finding \( f(5) \)}}
    The function \( f \) is defined by:
    \[
    f(x) = ab^{x/n},
    \]
    where \( a, b, \) and \( n \) are constants, and \( b, n \) are integers.  

    \vspace{0.3cm}
    If \( f(2) = 6 \) and \( f(4) = 150 \), what is the value of \( f(5) \)?
\end{frame}

% --- Answer Explanation 8 ---
\begin{frame}{\textbf{Answer Explanation: Finding \( f(5) \)}}
    \textbf{Correct Answer: \textcolor{novelPurple}{\( \mathbf{750} \)}}  

    \vspace{0.3cm}
    \textbf{Solution:}
    \begin{itemize}
        \item Given:
        \[
        f(2) = 6 \implies a \cdot b^{2/n} = 6.
        \]
        \[
        f(4) = 150 \implies a \cdot b^{4/n} = 150.
        \]
        \item Divide equations:
        \[
        \frac{b^{4/n}}{b^{2/n}} = b^{2/n}, \quad b^{2/n} = 25 \implies b^{1/n} = 5.
        \]
        \item Solve for \( f(5) \):
        \[
        f(5) = a \cdot b^{5/n} = \frac{6}{25} \cdot (5^5) = \frac{6}{25} \times 3125.
        \]
    \end{itemize}
    \textbf{Final Answer:} \( \boxed{750} \).
\end{frame}

\vspace{1cm}

\begin{frame}{\textbf{Question: Positive Difference Between \( a \) and \( b \)}}
    The quadratic equation graph of \( y = f(x) \) in the \( xy \)-plane has:
    \begin{itemize}
        \item A vertex at \( (4, -7) \).
        \item Contains the point \( (3, -9) \).
        \item Has a \( y \)-intercept at \( (0, a) \).
        \item The graph of \( y = 6 \cdot f(x) \) has a \( y \)-intercept at \( (0, b) \).
    \end{itemize}
    What is the positive difference between \( a \) and \( b \)?
\end{frame}

% --- Answer Explanation 9 ---
\begin{frame}{\textbf{Answer Explanation}}
    \textbf{Correct Answer: \textcolor{novelPurple}{\( \mathbf{195} \)}}  

    \vspace{0.3cm}
    \textbf{Solution:}
    \begin{itemize}
        \item Equation of the parabola:
        \[
        f(x) = k(x - 4)^2 - 7.
        \]
        \item Substitute \( (3, -9) \) to find \( k \):
        \[
        -9 = k(3 - 4)^2 - 7 \implies k = -2.
        \]
        \item Thus, the equation becomes:
        \[
        f(x) = -2(x - 4)^2 - 7.
        \]
        \item Find \( a \):
        \[
        a = f(0) = -39.
        \]
        \item Find \( b \) for \( y = 6 \cdot f(x) \):
        \[
        b = 6 \cdot (-39) = -234.
        \]

    \end{itemize}
\end{frame}


\begin{frame}{Answer Explanation}

    \textbf{Correct Answer: \textcolor{novelPurple}{\( \mathbf{195} \)}}  

    \vspace{0.3cm}
    \textbf{Solution:}
 \begin{itemize}
             \item Compute the positive difference:
        \[
        |a - b| = |234 - 39| = \boxed{195}.
        \]
 \end{itemize}
    
\end{frame}

% --- Question 10 ---
\begin{frame}{\textbf{Question: Finding \( a + b \)}}
    The function \( f \) is defined by:
    \[
    f(x) = ax^2 + bx + c,
    \]
    where \( a, b, \) and \( c \) are constants.  

    \vspace{0.3cm}
    The graph of \( y = f(x) \) in the \( xy \)-plane passes through the points \( (9,0) \) and \( (-2,0) \).  

    \vspace{0.3cm}
    If \( a \) is an integer greater than 1, which of the following could be the value of \( a + b \)?

    \vspace{0.5cm}
    \textbf{Options:}
    \begin{enumerate}
        \item[A] 8
        \item[B] 7
        \item[C] -6
        \item[D] -12
    \end{enumerate}
\end{frame}

% --- Answer Explanation 10 ---
\begin{frame}{\textbf{Answer Explanation: Finding \( a + b \)}}
    \textbf{Correct Answer: \textcolor{novelPurple}{\( \mathbf{-12} \)}}  

    \vspace{0.3cm}
    \textbf{Solution:}
    \begin{itemize}
        \item Factorize the quadratic equation:
        \[
        f(x) = a(x - 9)(x + 2) = a(x^2 - 7x - 18).
        \]
        \item From this equation:
        \[
        b = -7a, \quad c = -18a.
        \]
        \item For \( a \geq 2 \), calculate:
        \[
        a + b = a - 7a = -6a.
        \]
        \item Possible values:
        \[
        a + b = -12 \quad \text{when } a = 2.
        \]
    \end{itemize}
    \textbf{Final Answer:} \( \boxed{-12} \).
\end{frame}

% --- Question 11 ---
\begin{frame}{\textbf{Question: Finding a Factor of \( x + 2b \)}}
    Which of the following expressions has a factor of \( x + 2b \), where \( b \) is a positive integer constant?

    \vspace{0.3cm}
    \[
    f(x) =
    \begin{cases} 
    A: 3x^2 + 7x + 14b, \\ 
    B: 3x^2 + 22x + 14b, \\ 
    C: 3x^2 + 30x + 14b, \\ 
    D: 3x^2 + 37x + 14b.
    \end{cases}
    \]
\end{frame}

% --- Answer Explanation 11 ---
\begin{frame}{\textbf{Answer Explanation: Finding a Factor of \( x + 2b \)}}
    \textbf{Correct Answer: \textcolor{novelPurple}{\( \mathbf{D} \)}}  

    \vspace{0.3cm}
    \textbf{Solution:}
    \begin{itemize}
        \item By the factor theorem:
        \[
        \text{If } (x - c) \text{ is a factor of } f(x), \text{ then } f(c) = 0.
        \]
        \item Let \( x + 2b \) be a factor of \( f(x) \). Substitute \( x = -2b \):
        \[
        f(-2b) = 0.
        \]
        \item Calculate for each option, and we find that when \( b = 5 \), only:
        \[
        f(x) = 3x^2 + 37x + 14b
        \]
        satisfies \( f(-2b) = 0 \).
    \end{itemize}
    \textbf{Final Answer:} \( \boxed{D} \).
\end{frame}








\end{document}
